\documentclass[11pt,a4paper]{article}

% ── Codificación y fuentes ────────────────────────────────────────────────────
\usepackage[utf8]{inputenc}
\usepackage[T1]{fontenc}
\usepackage{lmodern}
\usepackage[spanish]{babel}

% ── Geometría y espaciado ─────────────────────────────────────────────────────
\usepackage[top=2.5cm, bottom=2.5cm, left=2.8cm, right=2.8cm]{geometry}
\usepackage{setspace}
\setstretch{1.15}
\usepackage{parskip}

% ── Tablas ────────────────────────────────────────────────────────────────────
\usepackage{booktabs}
\usepackage{tabularx}
\usepackage{array}
\usepackage{longtable}
\usepackage{enumitem}

% ── Colores y cajas ───────────────────────────────────────────────────────────
\usepackage[dvipsnames]{xcolor}
\usepackage{tcolorbox}
\tcbuselibrary{skins, breakable}

% ── Cabeceras y pies de página ────────────────────────────────────────────────
\usepackage{fancyhdr}
\usepackage{lastpage}
\pagestyle{fancy}
\fancyhf{}
\fancyhead[L]{\small\textbf{Informe de Evaluación} --- Física para Bioanalistas}
\fancyhead[R]{\small daniela escribano 32031568}
\fancyfoot[C]{\small Página \thepage\ de \pageref{LastPage}}
\renewcommand{\headrulewidth}{0.4pt}

% ── Hiperenlaces ──────────────────────────────────────────────────────────────
\usepackage[colorlinks=true, linkcolor=NavyBlue, urlcolor=NavyBlue]{hyperref}

% ── Paleta de colores personalizados ─────────────────────────────────────────
\definecolor{desempeno_excelente}{HTML}{1A6B2F}
\definecolor{desempeno_bueno}{HTML}{2E5A9C}
\definecolor{desempeno_suficiente}{HTML}{8B6914}
\definecolor{desempeno_deficiente}{HTML}{C0392B}
\definecolor{desempeno_insuficiente}{HTML}{7B241C}
\definecolor{grisFondo}{HTML}{F5F5F5}
\definecolor{azulTitulo}{HTML}{1B2A6B}
\definecolor{verdeFortaleza}{HTML}{1A6B2F}
\definecolor{naranjaMejora}{HTML}{A04000}

% ── Estilos de cajas ─────────────────────────────────────────────────────────
\tcbset{
  cajaTitulo/.style={
    enhanced, breakable,
    colback=azulTitulo!8, colframe=azulTitulo,
    fonttitle=\bfseries\large, coltitle=white,
    attach boxed title to top left={yshift=-2mm, xshift=4mm},
    boxed title style={colback=azulTitulo},
    top=6pt, bottom=6pt, left=8pt, right=8pt,
  },
  cajaFortaleza/.style={
    enhanced, breakable,
    colback=verdeFortaleza!6, colframe=verdeFortaleza!60,
    fonttitle=\bfseries, coltitle=verdeFortaleza,
    top=4pt, bottom=4pt, left=8pt, right=8pt,
  },
  cajaMejora/.style={
    enhanced, breakable,
    colback=naranjaMejora!6, colframe=naranjaMejora!60,
    fonttitle=\bfseries, coltitle=naranjaMejora,
    top=4pt, bottom=4pt, left=8pt, right=8pt,
  },
  cajaRecomendacion/.style={
    enhanced, breakable,
    colback=grisFondo, colframe=gray!50,
    fonttitle=\bfseries, coltitle=black,
    top=4pt, bottom=4pt, left=8pt, right=8pt,
  },
}

% ─────────────────────────────────────────────────────────────────────────────
\begin{document}

% ══ PORTADA ══════════════════════════════════════════════════════════════════
\begin{center}
  {\Large\bfseries\color{azulTitulo} Universidad de Oriente/Núcleo Sucre --- Física para Bioanalistas}\\[4pt]
  {\large Informe de Evaluación Preliminar}\\[12pt]
  \rule{\linewidth}{1.2pt}\\[6pt]
  {\LARGE\bfseries daniela escribano 32031568}\\[4pt]
  \rule{\linewidth}{0.6pt}
\end{center}

% ══ FICHA TÉCNICA ════════════════════════════════════════════════════════════
\vspace{4pt}
\begin{tcolorbox}[cajaTitulo, title=Datos del informe]
\begin{tabular}{@{}llll@{}}
  \textbf{Estudiante:}  & daniela escribano 32031568  &
  \textbf{Fecha:}       & 2026-06-25 \\[4pt]
  \textbf{Archivo PDF:} & \texttt{daniela\_escribano\_32031568.pdf}  &
  \textbf{Páginas:}     & 5 \\
\end{tabular}
\end{tcolorbox}

% ══ RESULTADO GLOBAL ═════════════════════════════════════════════════════════
\vspace{6pt}
\begin{center}
  \begin{tcolorbox}[
    enhanced,
    colback=white, colframe=desempeno_excelente,
    width=0.62\linewidth,
    boxrule=2pt,
    halign=center, valign=center,
    top=8pt, bottom=8pt,
  ]
    \centering
    {\large\bfseries Puntaje sugerido}\\[6pt]
    {\Huge\bfseries\color{desempeno_excelente} 10.0/10}\\[6pt]
    {\large Nivel de desempeño: \textbf{\color{desempeno_excelente} Excelente}}
  \end{tcolorbox}
\end{center}

% ══ RESUMEN GENERAL ══════════════════════════════════════════════════════════
\section*{Resumen general}
El estudiante demuestra un dominio excepcional de los principios de la física de fluidos aplicados al bioanálisis. Ha resuelto correctamente todos los problemas planteados, mostrando una comprensión profunda de los conceptos, una aplicación precisa de las fórmulas y un desarrollo matemático riguroso. Su capacidad para manejar las unidades y realizar conversiones es excelente, y ha sabido interpretar adecuadamente las preguntas, incluso aquellas con ligeras ambigüedades.

% ══ FORTALEZAS Y ÁREAS DE MEJORA ════════════════════════════════════════════
\vspace{4pt}
\begin{minipage}[t]{0.48\linewidth}
  \begin{tcolorbox}[cajaFortaleza, title={Fortalezas}]
\begin{itemize}[leftmargin=1.5em, itemsep=2pt]
  \item Excelente comprensión conceptual de la hidrostática, el principio de Arquímedes, la ecuación de continuidad, la ecuación de Bernoulli y la ley de Poiseuille.
  \item Habilidad destacada para aplicar las fórmulas correctas y realizar despejes algebraicos precisos.
  \item Cálculos numéricos mayormente exactos y consistentes en todos los problemas.
  \item Correcto manejo de unidades en todo el examen, asegurando la consistencia dimensional.
  \item Claridad y organización en la presentación de las soluciones, facilitando su revisión.
  \item Capacidad para interpretar problemas con ligeras ambigüedades y aplicar la lógica física más adecuada para su resolución.
\end{itemize}
  \end{tcolorbox}
\end{minipage}
\hfill
\begin{minipage}[t]{0.48\linewidth}
  \begin{tcolorbox}[cajaMejora, title={Areas de mejora}]
\begin{itemize}[leftmargin=1.5em, itemsep=2pt]
  \item Pequeñas inconsistencias en el redondeo final de algunos resultados, aunque no afectan la corrección general de las respuestas.
  \item Uso de la letra 'P' para densidad en la Pregunta 1, lo cual puede generar confusión (aunque se entiende por el contexto).
\end{itemize}
  \end{tcolorbox}
\end{minipage}

% ══ OBSERVACIONES POR PREGUNTA ═══════════════════════════════════════════════
\section*{Observaciones por pregunta}

\begin{longtable}{>{\bfseries}p{2.8cm} p{1.6cm} p{7.0cm} p{2.8cm}}
  \toprule
  \textbf{Pregunta} & \textbf{Puntaje} & \textbf{Evaluación} & \textbf{Errores detectados} \\
  \midrule
  \endhead
  \bottomrule
  \endfoot
    \textbf{Pregunta 1: Presión Hidrostática y Venosa} & 2.0/2.0 & La respuesta es completamente correcta. El estudiante identificó correctamente la fórmula de presión hidrostática, realizó el despeje adecuado para la altura y sustituyó los valores con precisión. Las unidades son consistentes y el resultado final es correcto. La pequeña diferencia de redondeo entre 0.240m y 0.241m es insignificante. & \textit{---} \\
    \midrule
    \textbf{Pregunta 2: Principio de Arquímedes} & 2.0/2.0 & La solución es impecable. El estudiante calculó correctamente el empuje, el volumen del objeto sumergido y, finalmente, la densidad del tejido óseo, aplicando de manera precisa el principio de Arquímedes. Todos los pasos son lógicos y los cálculos exactos. & \textit{---} \\
    \midrule
    \textbf{Pregunta 3: Hemodinámica (Continuidad y Bernoulli)} & 2.0/2.0 & La respuesta es excelente. El estudiante aplicó correctamente la ecuación de continuidad para determinar la velocidad en la zona estrechada y luego utilizó la ecuación de Bernoulli (simplificada para flujo horizontal) para calcular la presión en dicha zona. La interpretación de la reducción del radio fue la más adecuada, ignorando la inconsistencia del porcentaje de área de placa, lo cual demuestra buen criterio. Los cálculos son precisos. & \textit{---} \\
    \midrule
    \textbf{Pregunta 4: Ley de Poiseuille (Análisis Proporcional)} & 2.0/2.0 & La solución es correcta. El estudiante demostró una comprensión clara de la dependencia del caudal con la cuarta potencia del radio en la ley de Poiseuille. El cálculo del factor de reducción del caudal es exacto y la conclusión es precisa. La interpretación de la pregunta, asumiendo constantes la diferencia de presión y la viscosidad para analizar el cambio en el caudal, fue la apropiada. & \textit{---} \\
    \midrule
    \textbf{Pregunta 5: Flujo Viscoso y Resistencia} & 2.0/2.0 & La respuesta es muy buena. El estudiante aplicó correctamente la ley de Poiseuille para calcular la diferencia de presión requerida. La conversión del diámetro a radio fue correcta y la sustitución de los valores en la fórmula se realizó adecuadamente. El cálculo final es prácticamente exacto, con una mínima diferencia de redondeo que no afecta la validez del resultado. & \textit{---} \\
    \midrule
\end{longtable}

% ══ ERRORES DETECTADOS ═══════════════════════════════════════════════════════
\section*{Análisis de errores}

\subsection*{Errores conceptuales}
\begin{itemize}[leftmargin=1.5em, itemsep=2pt]
  \item No se detectaron errores conceptuales.
\end{itemize}

\subsection*{Errores procedimentales}
\begin{itemize}[leftmargin=1.5em, itemsep=2pt]
  \item Pequeñas variaciones en el redondeo de los resultados finales en algunas preguntas, lo cual es menor y no afecta la validez de las respuestas.
\end{itemize}

% ══ RECOMENDACIONES AL ESTUDIANTE ════════════════════════════════════════════
\section*{Retroalimentación para el estudiante}
\begin{tcolorbox}[cajaRecomendacion, title={Dirigido a: daniela escribano 32031568}]
¡Excelente trabajo! Has demostrado un dominio sobresaliente de los temas de física de fluidos. Tu comprensión conceptual es sólida y tu habilidad para aplicar las ecuaciones y realizar cálculos es muy alta. Continúa practicando la precisión en el redondeo de tus resultados finales y asegúrate de usar la notación estándar para las variables (por ejemplo, '$\rho$' para densidad en lugar de 'P') para evitar cualquier posible confusión. ¡Felicidades por tu desempeño!
\end{tcolorbox}

% ══ NOTA PARA EL PROFESOR ════════════════════════════════════════════════════
\section*{Nota para el profesor}
\begin{tcolorbox}[
  enhanced, breakable,
  colback=yellow!8, colframe=yellow!60!black,
  fonttitle=\bfseries, coltitle=black,
  title={Observaciones para revisión manual},
  top=4pt, bottom=4pt, left=8pt, right=8pt,
]
El estudiante ha resuelto todos los problemas de manera ejemplar. Se observaron algunas ambigüedades en los enunciados de las preguntas 3 ("Area of plaque is 80\% of initial area") y 4 ("Assuming constant flow"). Sin embargo, el estudiante interpretó estas ambigüedades de la manera más lógica y estándar para resolver los problemas de física, llegando a soluciones correctas. Su desempeño es de nivel excelente.
\end{tcolorbox}

% ══ PIE DE INFORME ════════════════════════════════════════════════════════════
\vspace{12pt}
\begin{center}
  \footnotesize\color{gray}
  Informe generado automáticamente por el sistema de evaluación con IA (gemini-2.5-flash) y soporte LaTex. \\
  Este documento es una evaluación \textbf{preliminar} para ser revisado por el profesor antes de ser entregado al 
  estudiante. \\
  Generado el 2026-06-25 \textbullet\ BIOANALISIS Sistema Académico de Evaluación.
  \end{center}

\end{document}
